\begin{ouverture}[courbe]{9}{Relier des grandeurs et lire des graphiques}{Partie 2 · Algèbre et modèles}
{Tableaux · Formules · Courbes · Interprétation}
\objectif{Traduire une dépendance par une formule et un tableau.}
\objectif{Lire dans les deux sens sur un graphique.}
\objectif{Distinguer proportionnalité et autre relation.}
\prerequis{Repérage dans le plan et calcul littéral.}
\end{ouverture}
\pagestyle{otMain}\essentiel\fontsize{11}{15}\selectfont
\phantomsection\label{cours-9-1}
\essSection{1}{Trois représentations}
\begin{definition}Une grandeur peut dépendre d’une autre. Une formule donne une règle de calcul, un tableau donne des couples de valeurs, un graphique les représente. En Quatrième, on travaille ces correspondances sans formaliser les notations de fonctions.\end{definition}
\begin{exemple}Une location coûte 6 euros fixes puis 2 euros par heure. Pour une durée $t$ en heures, le prix est $P=6+2t$. Pour $t=0,1,2,3$, les prix sont $6,8,10,12$ euros. Ce n’est pas une proportionnalité car le prix n’est pas nul pour $t=0$.\end{exemple}
\label{fig-cours-9-1-cours-location}\figureGeo{cours-location}

\begin{savoirfaire}[Construire]Nommer et graduer les axes avec leurs unités. Placer la grandeur choisie sur l’axe horizontal. Calculer et placer plusieurs couples. Relier seulement si le contexte définit les valeurs intermédiaires, et avec une forme justifiée par la règle ou explicitée dans l’énoncé.\end{savoirfaire}
\clearpage
\phantomsection\label{cours-9-2}
\essSection{2}{Lire dans les deux sens}
\begin{savoirfaire}[De l’horizontale vers la verticale]Repérer une valeur sur l’axe horizontal, rejoindre la courbe verticalement, puis lire la valeur correspondante sur l’axe vertical. Dans l’autre sens, partir d’une valeur verticale et suivre une horizontale jusqu’à toutes les intersections avec la courbe.\end{savoirfaire}
\begin{exemple}Une température de four peut être atteinte à la montée et à la descente : la lecture inverse peut donner plusieurs durées. Une hauteur non atteinte n’a pas de durée correspondante dans l’intervalle représenté.\end{exemple}
\begin{propriete}Une proportionnalité se représente par une droite passant par l’origine. Une droite qui ne passe pas par l’origine représente une autre relation. Une courbe peut représenter une dépendance non proportionnelle, par exemple l’aire $A=c^2$ d’un carré de côté $c$.\end{propriete}
\begin{attention}Lire d’abord les graduations et les unités : une graduation peut valoir 2 minutes ou 20 degrés. Une lecture graphique est souvent approchée. Le graphique ne dit rien au-delà de la zone et des hypothèses précisées ; prolonger une tendance demande une justification.\end{attention}

\clearpage\rubrique{Découverte}
\begin{blocLivre}{ex-9-1}\begin{exercice}[1]{Une formule de tarif}
Une location coûte 5 euros fixes et 3 euros par heure. Exprimer le prix $P$ pour $t$ heures.
\end{exercice}\end{blocLivre}
\begin{blocLivre}{ex-9-2}\begin{exercice}[1]{Compléter les valeurs}
Pour $P=5+3t$, calculer les prix pour $t=0$, $t=2$ et $t=4$.
\end{exercice}\end{blocLivre}
\begin{blocLivre}{ex-9-3}\begin{exercice}[1]{Lire un couple}
Sur le graphique de température, lire la température à 2 min.
\label{fig-ex-9-3}\figureGeo{four}
\end{exercice}\end{blocLivre}
\begin{blocLivre}{ex-9-4}\begin{exercice}[1]{Lire une durée}
Sur le même graphique, à quel instant la température atteint-elle $100\,^\circ$C pour la première fois ?
\label{fig-ex-9-4}\figureGeo{four}
\end{exercice}\end{blocLivre}
\begin{blocLivre}{ex-9-5}\begin{exercice}[1]{Reconnaître une relation}
Le graphique représente-t-il une situation proportionnelle ? Justifier.
\label{fig-ex-9-5}\figureGeo{graph-non-proportion}
\end{exercice}\end{blocLivre}

\rubrique{Entraînement}
\begin{blocLivre}{ex-9-6}\begin{exercice}[2]{Construire un graphique}
On vend une corde à 1,5 euro par mètre, pour des longueurs de 0 à 6 m. Compléter un tableau pour 0, 2, 4 et 6 m, puis tracer la représentation.
\label{fig-ex-9-6}\figureGeo{graph-corde}
\end{exercice}\end{blocLivre}
\begin{blocLivre}{ex-9-7}\begin{exercice}[2]{Une lecture inverse double}
Sur le graphique de température, donner tous les instants où la température vaut $60\,^\circ$C.
\label{fig-ex-9-7}\figureGeo{four}
\end{exercice}\end{blocLivre}
\begin{blocLivre}{ex-9-8}\begin{exercice}[2]{Aire d’un carré}
Exprimer l’aire $A$ d’un carré de côté $c$ cm. Calculer pour $c=1,2,3,4$. Y a-t-il proportionnalité ?
\label{fig-ex-9-8}\figureGeo{graph-carre}
\end{exercice}\end{blocLivre}
\begin{blocLivre}{ex-9-9}\begin{exercice}[2]{Une formule à retrouver}
Un tableau donne les couples $(0;4)$, $(1;7)$, $(2;10)$ et $(3;13)$. On sait que le modèle est un montant fixe plus 3 fois la première valeur. Écrire la formule.
\end{exercice}\end{blocLivre}
\begin{blocLivre}{ex-9-10}\begin{exercice}[2]{Un tableur}
Dans un tableur, A2 contient une durée en heures. Quelle formule saisir en B2 pour le prix $P=5+3t$, puis recopier vers le bas ?
\label{fig-ex-9-10}\figureGeo{tableur-prix}
\end{exercice}\end{blocLivre}

\rubrique{Problèmes}
\begin{blocLivre}{ex-9-11}\begin{exercice}[3]{Retirer des carrés}
Dans un rectangle de 18 cm sur 10 cm, on retire aux quatre coins des carrés de côté $x$ cm, avec $0<x<5$. Exprimer l’aire restante et la calculer pour $x=2$.
\label{fig-ex-9-11}\figureGeo{coins-retires}
\end{exercice}\end{blocLivre}
\begin{blocLivre}{ex-9-12}\begin{exercice}[3]{Deux courbes de prix}
Les graphiques correspondent aux prix A $=4+2t$ et B $=4t$, pour $0\le t\le4$ h. Lire l’heure d’égalité et vérifier par une équation.
\label{fig-ex-9-12}\figureGeo{graph-tarifs}
\end{exercice}\end{blocLivre}
\begin{blocLivre}{ex-9-13}\begin{exercice}[3]{Une réserve qui diminue}
Une réserve contient 120 L ; on retire 8 L par minute à débit constant. Donner une formule pour le volume après $t$ minutes et la durée de validité avant épuisement.
\label{fig-ex-9-13}\figureGeo{graph-reserve}
\end{exercice}\end{blocLivre}
\begin{blocLivre}{ex-9-14}\begin{exercice}[3]{Connaître les limites}
Un graphique ne montre que les ventes de janvier à juin. Peut-on affirmer que les ventes de décembre seront égales au double de celles de juin si le tracé monte ?
\end{exercice}\end{blocLivre}
\begin{blocLivre}{ex-9-15}\begin{exercice}[3]{Comparer deux dépendances}
Un carré de côté $c$ cm a un périmètre $P=4c$ et une aire $A=c^2$. Pour quels côtés positifs leur valeur numérique est-elle la même ? Chercher un exemple, puis justifier sans résoudre d’équation du second degré.
\label{fig-ex-9-15}\figureGeo{carre-4}
\end{exercice}\end{blocLivre}

\clearpage
